We study the effect of a spatially varying diffusivity $D(x)$ on the mean first
passage time (MFPT) to a partially absorbing target.  The motivation is foraging in
crowded environments: local crowder density modulates the effective diffusion
coefficient, raising the question of whether a non-uniform crowding profile can
speed up or slow down search.  We work in 1D on the ring $[-\pi,\pi]$ with periodic
boundary conditions and a partially absorbing trap of strength $\kappa$ at $x=0$,
modelled by a $\delta$-sink in the Fokker--Planck equation.  The lattice random walk
underlying the model is derived, and Fick's form
$\partial_t p=\partial_x[D(x)\partial_x p]$ is justified by the symmetric exclusion
process.  From the backward (adjoint) equation we obtain the MFPT in the closed
form
\[
\langle T\rangle
\;=\;\frac{2\pi}{\kappa}\;+\;\frac{1}{\pi}\!\int_0^\pi\!\frac{(\pi-y)^2}{D(y)}\,dy,
\]
which decomposes into an additive dwell plus travel time.
For the piecewise-linear (tent) profile $D(x)=D_0+D_1(1-|x|/\pi)$ we evaluate the
travel integral explicitly and recover a logarithmic closed form.  The analytical
result is verified in three independent ways: (i) a finite-difference solution of
the backward equation, with empirical second-order convergence; (ii) a
Crank--Nicolson finite-difference solution of the forward Fokker--Planck equation,
from whose survival probability the MFPT is recovered by time integration; and
(iii) Monte Carlo simulation of the It\^o SDE $dx=D'(x)\,dt+\sqrt{2D(x)}\,dW$ with
a Poisson trap rule.  All three numerical methods agree with the analytical
formula within their expected error.  Finally, we examine three microscopic
stories for $\kappa$ as a function of $D(0)$ and show that the biologically
motivated $\kappa\propto 1/D(0)$ story is the only one producing a non-trivial
optimal crowding level.
